\section{系统级复盘：GLOM 到底组合了什么}

走到这里，可以把 GLOM 压成一个 recurrent structured inference loop：低层 evidence 向上提出 parent hypothesis，高层 context 向下预测 part，same-level attention 让相似 hypothesis 形成 islands，previous state 提供 inertia，masked reconstruction 与 consensus distillation 训练这些函数，location-conditioned decoder 处理 shared whole 到不同 parts 的生成。

\lecturefigure{slide-32-summary.jpg}{课堂总结：用 Transformer、contrastive learning 与 neural field 组合 GLOM。}{Stanford Online 官方视频 00:50:08--00:52:20。}

\subsection{机制清单与失败模式}

\begin{table}[H]
\centering
\small
\caption{GLOM 的机制、解决的问题与尚缺证据。}
\begin{tabularx}{\textwidth}{L{0.20\textwidth}L{0.30\textwidth}>{\raggedright\arraybackslash}X}
\toprule
机制 & 试图解决的问题 & 尚缺的验证 \\
\midrule
Dynamic embeddings & 固定硬件表示输入特定 parse & 与 ground-truth hierarchy 的 alignment \\
Islands of agreement & 用相同 activity 表示同一 node & 防 collapse、跨实例分离、边界稳定性 \\
Bottom-up/top-down maps & part-whole identity/pose consistency & coordinate equivariance 与 occlusion robustness \\
Local similarity attention & 发现 spatial coherence region & attractor、error propagation 与 convergence \\
Hough-style parent voting & 避免 pairwise relation routing explosion & 对复杂 scene、多实例与 deformable object 的效果 \\
Consensus distillation & 让 local predictors 共享结构信息 & teacher drift、错误共识与 optimization stability \\
Sparse upper levels & 扩大 range 而控制 cost & 实际 FLOPs、memory、latency 与 communication \\
Neural field decoder & 相同 whole 在不同位置预测不同 parts & 与 islands 联合训练后的生成质量 \\
\bottomrule
\end{tabularx}
\end{table}

\subsection{读论文时应追问的实验}

若今天把 GLOM 从 design document 推向系统论文，至少需要以下实验闭环：

\begin{enumerate}
  \item \textbf{Representation readout}：定义如何从连续 embeddings 提取 nodes、edges、identity 与 pose；
  \item \textbf{Causal use}：干预某个 island 后，验证对应 part/whole prediction 是否按结构变化；
  \item \textbf{Convergence}：测 settling steps、fixed-point stability、oscillation 与 early stopping；
  \item \textbf{Anti-collapse}：比较 local attention、negative pairs、variance regularization 与 consensus weighting；
  \item \textbf{Generalization}：测试新 viewpoint、新 part composition、occlusion、multiple instances 与 video tracking；
  \item \textbf{Systems accounting}：报告 state memory、attention edges、BPTT activation、wall-clock latency 与 energy；
  \item \textbf{Ablation}：分别移除 top-down、bottom-up、persistence、lateral attention、location condition 与 distillation。
\end{enumerate}

\begin{importantbox}{一个可证伪的 GLOM 主张}
若 islands 真正承担 parse-tree node，移除或扰动某个 parent-level island，应系统性改变其 constituent parts 的 top-down prediction，而不只是降低全局分类分数。这个 intervention 比“可视化看起来像 object”更接近机制证据。
\end{importantbox}

\subsection{本章小结}

GLOM 的强项是把 representation、inference、learning 与 coordinate generation 放进同一设计语言；弱项是几乎每条关键链路都缺少端到端实验。正因为证据不完整，它特别适合作为 architecture thinking exercise：强迫读者把“对象结构”从直觉词汇改写成 state、message、loss 与 verification contract。

